\section{Conclus\~ao e Trabalhos Futuros}
\label{sec:conclusao}

A pergunta deste trabalho foi se as falas anteriores de uma pessoa, condensadas num perfil, ajudam
um LLM a reconhecer a opinião que a imprensa lhe atribuiu numa audiência posterior. Para isso, o
pipeline liga as opiniões publicadas sobre audiências públicas à fala de quem as emitiu e condensa a
fala de cada ator em um perfil escrito somente a partir dela. Esse perfil é usado para simular como a
pessoa se posicionaria, com a base da estimativa declarada.

A UDV localiza evidência candidata para 2.105 das 2.203 opiniões, mas, num controle com as 183
opiniões de treino que contêm citação confiável, a sentença de maior cosseno coincide com a passagem
citada em 61,2\% dos casos. A precisão dos vínculos por similaridade semântica ainda precisa ser
medida pela validação humana. <Síntese dos resultados da simulação.>

Ficam em aberto uma condição só com os trechos recuperados e a comparação entre perfis gerados por
modelos e prompts diferentes. A resposta simulada é uma estimativa condicionada às falas públicas da
pessoa e não equivale a uma declaração dela. Por isso, deve ser apresentada com a base da
estimativa, nunca como citação, e não serve para inferir intenção política.
